\section{Related Work}
\label{sec:related-work}

% Structure by ASPECT, not as a list of papers (course writing recs).
% Stress what makes answer-side-only translation different.

\subsection{Commonsense question answering}
% CommonsenseQA \citep{talmor-etal-2019-commonsenseqa}: construction, ConceptNet
% basis, the validation/test split situation (test labels withheld).

\subsection{Multilingual commonsense reasoning}
% X-CSQA \citep{lin-etal-2021-common} — full-item translation across many languages.
% mCSQA \citep{sakai-etal-2024-mcsqa} — unified LM+human multilingual construction.
% Position our work: these translate the WHOLE item or build new items; we hold the
% question fixed in English and translate ONLY the answer choices.

\subsection{Code-switching and cross-lingual grounding}
% Prior code-switching work tends to be question-side. Frame answer-side-only
% translation as an underexplored diagnostic for cross-lingual concept grounding.

% TODO: prose; cite conference versions verified against the ACL Anthology.
